\chapter{System Requirements and Architecture}
\label{chap:architecture-requirements}
% ============================================================

\section{Functional Requirements}

Table~\ref{tab:functional-requirements} lists the functional requirements for the complete interaction and execution stack. They cover natural-language ingress, dialogue routing, symbolic grounding, planning, deterministic execution, feedback, and observability. Their implementation owners are introduced in Section~\ref{sec:high-level-ros-architecture} and detailed in Chapter~\ref{chap:implementation}.

\begin{table}[H] \centering \small \begin{tabular}{p{0.14\linewidth}p{0.78\linewidth}} \toprule \textbf{ID} & \textbf{Requirement}\\ \midrule FR1 & The system shall accept natural-language user input through the dialogue stack.\\
FR2 & The system shall distinguish dialogue-only turns from execution-eligible turns.\\
FR3 & The system shall convert execution-eligible user turns into admitted planner requests before invoking the planner.\\
FR4 & The planner shall generate structured skill plans over an explicit skill registry.\\
FR5 & The orchestrator shall validate planner outputs before any robot skill is dispatched.\\
FR6 & The system shall publish execution feedback after relevant plan and step/skill events.\\
FR7 & The planner shall support asking for clarification, replanning and plan failure when execution feedback indicates that the current plan cannot continue safely.\\
FR8 & The system shall expose traces, logs, or runtime messages suitable for validation of the stack.\\
\bottomrule \end{tabular} \caption{Functional requirements for the proposed NAO interaction and execution stack.} \label{tab:functional-requirements} \end{table}

\section{Design Principles}

Four responsibility principles guide the implementation of Table~\ref{tab:functional-requirements}. They define the component roles before the corresponding ROS~2 package names are introduced.

\begin{enumerate}
  \item \textbf{Dialogue remains distinct from execution.} Greetings, conversation, and questions about the robot's known state remain non-executable unless the user explicitly requests an action.
  \item \textbf{The planner proposes structured skill plans.} It selects from the reviewed registry and uses grounded scene context, but does not execute skills, speak directly, or call robot-specific APIs.
  \item \textbf{The orchestrator remains the execution gate.} It admits execution requests, validates plans, dispatches skills, and publishes typed execution feedback.
  \item \textbf{Skills own fresh execution evidence.} Perception and execution skills query the knowledge base (KB) or robot state when required. Planning-time context cannot establish that a later action succeeded.
\end{enumerate}

% ============================================================
\section{High-Level ROS 2 Architecture}
\label{sec:high-level-ros-architecture}
% ============================================================

The responsibilities above are implemented by concrete ROS~2 packages. This section introduces the node names used in later chapters for contracts, source ownership, and runtime evidence. The architecture keeps dialogue, semantic routing, grounding, planning, deterministic execution, and feedback in separate components.

The main nodes and shared representations are:

\begin{itemize}
\item \code{dialogue_manager} acts as the central dialogue ownership component, maintaining conversational state, coordinating turn-taking, managing interaction flow, and serving as the integration point between user-facing dialogue and downstream planning services.

\item \code{chatbot_llm} is responsible for robot utterance generation, intent declaration and goal extraction during planner-relevant turns. It transforms natural language interactions into structured representations that depict the nature of the interaction. \textbf{Dialogue} and \textbf{knowledge query }turns remain conversational, \textbf{execution} turns trigger a planner request.

\item \code{planner_llm} performs proactive task planning over the available skill registry, normalized user intents, dialogue context, and current robot world state. The planner reasons over high-level objectives, decomposes user requests into executable actions, and generates candidate plans that align with both robot capabilities and current world state constraints.

\item \code{nao_orchestrator} provides deterministic plan validation and skill execution management. It verifies planner outputs against available skills and execution constraints, coordinates skill invocation, monitors execution progress, and emits structured execution feedback that can be consumed by both the dialogue and planning layers.

\item \code{grounded_context} is the compact KB-to-language JSON projection defined by Contract~3. It combines selected knowledge statements, the current scene summary, and tracked-person state into role-separated entities and bounded relations for the chatbot and planner.

\item Skill servers expose typed execution interfaces over lower-level ROS components and robot capabilities. They encapsulate perception, navigation, manipulation, motion, and communication actions. The orchestrator invokes these interfaces after plan validation, and their results provide fresh evidence for feedback, reporting, and verified KB effects.

\end{itemize}

Figure~\ref{fig:layered-architecture} shows how the packages are grouped into dialogue, grounding, planning, and execution layers. It also shows the intent, context, plan, and feedback paths that connect them.

% ============================================================
% High-level layered architecture
% ============================================================

\begin{figure}[H]
\centering
\resizebox{\textwidth}{!}{%
\begin{tikzpicture}[
  font=\normalsize,
  >=Latex,
  box/.style={
    draw,
    rounded corners,
    align=center,
    inner sep=6pt
  },
  topbox/.style={
    box,
    text width=6.1cm,
    minimum height=1.70cm
  },
  widebox/.style={
    box,
    text width=13.2cm,
    minimum height=1.30cm
  },
  control/.style={-{Latex[length=2.7mm]}, thick},
  support/.style={-{Latex[length=2.4mm]}, thick, dashed},
  label/.style={font=\scriptsize, fill=white, inner sep=1.3pt}
]

% Top row: dialogue and grounding are parallel inputs.
\node[topbox, fill=blue!7] (dialogue) at (-4.45,0)
{\textbf{Dialogue Layer}\\
\code{dialogue_manager} + \code{chatbot_llm}\\
turn lifecycle, speech ownership, route and intent};

\node[topbox, fill=orange!10] (grounding) at (4.45,0)
{\textbf{Grounding Layer}\\
\code{nao_scene_grounding}, \code{kb_skills}, KnowledgeCore\\
symbolic state, scene summaries, entity grounding};

% Main layers, with enough vertical spacing to avoid overlap.
\node[widebox, fill=green!8] (planning) at (0,-4.35)
{\textbf{Planning Layer}\\
\code{planner_llm}: structured skill-plan generation, feedback-driven goal-state handling, replanning decisions, dialogue-act intent};

\node[widebox, fill=purple!7] (execution) at (0,-7.10)
{\textbf{Execution Layer}\\
\code{nao_orchestrator} + skill servers: planner admission, validation, ordered dispatch, skill results, execution feedback};

% Grounding also supports chatbot interpretation.
\draw[support]
  (grounding.west) -- node[label, below]{scene / symbolic context}
  (dialogue.east);

% Parallel inputs into the planner.
\draw[control]
  (dialogue.south) -- node[label, left]{task metadata / planner request}
  ([xshift=-3.35cm]planning.north);

\draw[support]
  (grounding.south) -- node[label, right]{compact grounded context}
  ([xshift=3.35cm]planning.north);

% Planner to execution.
\draw[control]
  (planning.south) -- node[label, right]{validated structured plan}
  (execution.north);

% Feedback loop, kept close to avoid shrinking the whole figure.
\draw[control]
  ([yshift=0.10cm]execution.west) to[out=180,in=180,looseness=0.75]
  node[label, left]{execution feedback / goal-state handling}
  ([yshift=-0.10cm]planning.west);

% Evidence update loop, kept close to avoid widening the bounding box.
\draw[support]
  ([yshift=0.10cm]execution.east) to[out=0,in=0,looseness=0.75]
  node[label, right]{skill evidence / KB updates}
  ([yshift=-0.10cm]grounding.east);

\end{tikzpicture}%
}
\caption{High-level layered architecture. Dialogue interpretation and symbolic grounding provide parallel inputs to the planning layer. The chatbot layer also consumes symbolic scene context. The planner proposes structured skill plans, while the execution layer validates and dispatches skills before returning feedback for replanning or dialogue.}
\label{fig:layered-architecture}
\end{figure}

% ============================================================
% Runtime flow
% ============================================================

\subsection{Runtime Flow}

Figure~\ref{fig:runtime-sequence} expands the layered view into the normal execution sequence. The numbered stages distinguish user-turn ownership, planner admission, candidate plan generation, deterministic dispatch, fresh skill evidence, and the return path for replanning or user communication.

\begin{figure}[H]
\centering
\resizebox{0.90\textwidth}{!}{%
\begin{tikzpicture}[
  font=\small,
  >=Latex,
  node distance=0.58cm,
  box/.style={
    draw,
    rounded corners,
    align=center,
    text width=6.8cm,
    minimum height=0.82cm,
    fill=gray!5,
    inner sep=4pt
  },
  dialogue/.style={box,fill=blue!7},
  planning/.style={box,fill=green!8},
  execution/.style={box,fill=purple!7},
  evidence/.style={
    draw,
    rounded corners,
    align=center,
    text width=3.9cm,
    minimum height=0.82cm,
    fill=orange!10,
    inner sep=4pt
  },
  arr/.style={-{Latex[length=2.4mm]},line width=0.65pt},
  side/.style={-{Latex[length=2.2mm]},line width=0.55pt,dashed},
  edgelabel/.style={font=\scriptsize,fill=white,inner sep=1.2pt}
]

\node[box] (u) {1. User utterance};

\node[dialogue,below=0.62cm of u] (dm)
{2. \code{dialogue_manager}: receive dialogue turn and invoke chatbot};

\node[dialogue,below=0.62cm of dm] (c)
{3. \code{chatbot_llm}: derive intent, route, goal text, and compact context};

\node[execution,below=0.62cm of c] (g)
{4. \code{nao_orchestrator}:\\ \textbf{planner\_gate} admits execution request};

\node[planning,below=0.62cm of g] (p)
{5. \code{planner_llm}: generate structured skill plan};

\node[execution,below=0.62cm of p] (o)
{6. \code{nao_orchestrator}: dispatch executable steps after validation};

\node[execution,below=0.62cm of o] (s)
{7. Skills: execute action and refresh live evidence};

\node[execution,below=0.62cm of s] (f)
{8. \code{nao_orchestrator}: publish execution feedback};

\node[planning,below=0.62cm of f] (sup)
{9. \code{planner_llm}: process result, replan, or request communication};

\node[dialogue,below=0.62cm of sup] (d)
{10. \code{dialogue_manager}: realize spoken response};

\node[evidence,right=1.10cm of c] (kb1)
{Scene / KB\\compact projection};

\node[evidence,right=1.10cm of s] (kb2)
{Raw scene / KB state};

\draw[arr] (u) -- (dm);
\draw[arr] (dm) -- node[edgelabel,right]{chatbot turn call} (c);
\draw[arr] (c) -- node[edgelabel,right]{\code{/nao_orchestrator/planner_request}} (g);
\draw[arr] (g) -- node[edgelabel,right]{\code{/planner/request}} (p);
\draw[arr] (p) -- node[edgelabel,right]{sequence of steps/skills} (o);
\draw[arr] (o) -- node[edgelabel,right]{skill action/service call} (s);
\draw[arr] (s) -- node[edgelabel,right]{skill result} (f);
\draw[arr] (f) -- node[edgelabel,right]{\code{/planner/execution_feedback}} (sup);
\draw[arr] (sup) -- node[edgelabel,right]{\code{/planner/dialogue_act}} (d);

% Grounding/evidence support paths.
\draw[side] (kb1) -- node[edgelabel,above]{symbolic context} (c);
\draw[side] (s) -- (kb2);
\draw[side] (kb2) -- (s);

\end{tikzpicture}%
}
\caption{Runtime message and responsibility sequence from user input to spoken response. The colour coding follows Figure~\ref{fig:layered-architecture}: dialogue components are blue, grounding and evidence are orange, planning is green, and execution is purple.}
\label{fig:runtime-sequence}
\end{figure}

% ============================================================
